\pagenumbering{roman}
\setcounter{secnumdepth}{-1}

\addtocontents{toc}{\protect\setcounter{tocdepth}{-1}}

\setcounter{page}{0}

\chapter{Linear Algebra Done Right}

\thispagestyle{empty}

\begin{center}

Fourth Edition

\medskip

19 November 2023
\\\copyright\ 2024 Sheldon Axler

\vfill

{\huge
Sheldon Axler}

\vfill

Comments, corrections, and suggestions\\about this book are most welcome.\\Please send them to \PClink{linear@axler.net}.

\vfill

The print version of this book is published by Springer.

\bigskip

\end{center}

\CrCom


\vfill

{\changefontsize{16}

\[
F_n = \frac{1}{\sqrt{5}} \, \Biggl[ \paren[\Bigg]{ \frac{1 + \sqrt{5}}{2} }^{\!n}
-\paren[\Bigg]{ \frac{1 - \sqrt{5}}{2} }^{\!n} \Biggr]
\]}

\addtocontents{toc}{\protect\setcounter{tocdepth}{2}}

\setcounter{page}{4}

\chapter{About the Author}

Sheldon Axler received his undergraduate degree from Princeton University, followed by a PhD in mathematics from the University of California at Berkeley.

%Sheldon Axler was valedictorian of his high school class in Miami, Florida. He received his AB from Princeton University with highest honors, followed by a PhD in mathematics from the University of California at Berkeley.

As a postdoctoral Moore Instructor at MIT, Axler received a university-wide teaching award. He was then an assistant professor, associate professor, and professor at Michigan State University, where he received the first J. Sutherland Frame Teaching Award and the Distinguished Faculty Award.

Axler received the Lester R. Ford Award for expository writing from the Mathematical Association of America in 1996, for a paper that eventually expanded into this book. In addition to publishing numerous research papers, he is the author of six mathematics textbooks, ranging from freshman to graduate level. Previous editions of this book have been adopted as a textbook at over 375 universities and colleges and have been translated into three languages.

Axler has served as Editor-in-Chief of the \textit{Mathematical Intelligencer} and Associate Editor of the \textit{American Mathematical Monthly}. He has been a member of the Council of the American Mathematical Society and of the Board of Trustees of the Mathematical Sciences Research Institute. He has also served on the editorial board of Springer's series Undergraduate Texts in Mathematics, Graduate Texts in Mathematics, Universitext, and Springer Monographs in Mathematics.

Axler is a Fellow of the American Mathematical Society and has been a recipient of numerous grants from the National Science Foundation.

Axler joined San Francisco State University as chair of the Mathematics Department in 1997. He served as dean of the College of Science \& Engineering from 2002 to 2015, when he returned to a regular faculty appointment as a professor in the Mathematics Department.

\smallskip

\FigureHereCredit{artwork/SheldonMoonGGBridge.png}{0.39}{The author and his cat Moon.}{0in}{Carrie Heeter, Bishnu Sarangi} \label{0moon}\index{Moon}

\medskip

\noindent
\textit{Cover equation}: Formula for the $n^\nth$ Fibonacci number. Exercise \ref{Fibonacci} in Section~\ref{diagmat} uses linear algebra to derive this formula.

{\overfullrule=0in
\addtocontents{toc}{\protect\setcounter{tocdepth}{-1}}

\parfillskip=14.2bp plus1fil
%\renewcommand{\footfix}{\protect\footnotemark[1]}
\chapter{Contents}\tableofcontents}

\addtocontents{toc}{\protect\setcounter{tocdepth}{2}}

\endinput 
